\documentclass[10pt,a4paper]{amsart}
\usepackage[margin=19mm]{geometry}
\usepackage{amsmath,amssymb,mathtools}
\usepackage[scr=rsfs,cal=euler]{mathalfa}
\usepackage{xfrac}
\usepackage[unicode,hidelinks,bookmarks=false]{hyperref}
\newcommand{\ie}{i.e.\ }
\newcommand{\cf}{cf.\ }
\newcommand{\resp}{respectively\ }
\newcommand{\Subsection}{\subsection}
\providecommand{\nobreakdash}{\nobreak\hbox{-}}
\setlength{\emergencystretch}{2em}
\multlinegap0pt
\usepackage{amssymb}
\usepackage{enumerate}
\usepackage{float}
\usepackage{xcolor}
\usepackage{bm}
\usepackage{braket}
\usepackage[shortlabels]{enumitem}
\newtheorem{thm}{Theorem}
\DeclareMathOperator{\Id}{Id}
\title{Linear representation of groups associated to graphs}
\author{Alain Bretto, Alain Faisant, Neda Moradi, Mehrdad Nasernejad}
\date{Source: arXiv:2607.15287v1 (CC BY 4.0)}
\begin{document}
\maketitle

\noindent\emph{Source and licence.} Linear representation of groups associated to graphs. Version arXiv:2607.15287v1 (CC BY 4.0), distributed by arXiv under the Creative Commons Attribution 4.0 International licence, http://creativecommons.org/licenses/by/4.0/, as recorded in the arXiv licence metadata for this exact version and shown on the abstract page https://arxiv.org/abs/2607.15287v1. Full licence notice: \texttt{licenses/arxiv-2607.15287-CC-BY-4.0.txt}.\par\smallskip

\noindent\textit{Editorial scope.} Counted unit: the article's thm graphEquivModules (source0.tex lines 658--666). The article's complete original proof of it is delivered with it (source0.tex lines 668--713, one \texttt{proof} environment of 46 lines). No mathematics is written, repaired or re-derived: the statement and the proof are the article's own lines, in the article's own notation, and no retained line is substituted or re-flowed.  No further source-local context is needed: every cross-reference of every retained line resolves inside the two delivered spans.  Nothing was omitted from any retained span, so the package carries no omission record.  No retained line contains a citation, so the wrapper carries no bibliography and no bibliography entry.  No retained line contains a figure environment, an image-inclusion command or drawing code, so no image is delivered and the package declares no asset dependency.  Editorial formatting only, in addition: an AMS \texttt{amsart} preamble that restates the article's own declarations and macros used by the retained lines, namely the environments \texttt{thm} on separate counters, together with the standard packages those lines need.  The retained environments print in this document as Theorem 1 in order of appearance, whereas the article prints the same items with the numbering of its own full document.  The complete native source of the article as distributed in the arXiv e-print for this exact version is bundled unchanged as \texttt{sources/205372/source0.tex}.
\begin{thm} \label{graphEquivModules}
Let $R_m=\overline{\mathbb{F}}_p[\mathbb{E}_m]$, with  $\mathbb{E}_m=\mathbb{F}_2[\mathcal{E}_m], \; m \ge 2$.\\
 Let  $\Gamma_m=(V;\mathcal{E}_m, \alpha)$ and $\Gamma'_m=(V'; \mathcal{E}_m, \alpha')$ be two simple $m$-graphs. The following properties are equivalent:
\begin{enumerate}
\item[i)] the graph $\Gamma_m$ is $m$-isomorphic to  $\Gamma'_m$;
\item[ii)] the associated $R_m$-modules ${\displaystyle \mathbb{V}(\Gamma_m)}$ and ${\displaystyle \mathbb{V}(\Gamma'_m)}$ are isomorphic;
\item[iii)]  the representations $ \rho_{\Gamma_m}$ and $\rho_{\Gamma'_m}$  are equivalent.
\end{enumerate}
\end{thm}

\begin{proof}\mbox{}

\begin{itemize}
\item[-] $i)\Longrightarrow $ii) : let $(\varphi, \Id) : \Gamma_m \longrightarrow \Gamma'_m$ be an $m$-isomorphism: $\varphi \circ \alpha=\alpha'$; we have: 
$f(\alpha(e))=\alpha'(\Id(e))= \alpha'(e)$, consequently:
\begin{equation*}
\begin{array}{l}
\left((x\in \alpha(e)\Longleftrightarrow \varphi(x)\in \alpha'(e)\right)\Longleftrightarrow\\
\left(e\in \gamma(x)\Longleftrightarrow  e\in \gamma'(\varphi(x))\right)\Longleftrightarrow\\
 \left(\xi_{\gamma(x)}\oplus R_m\right)\simeq \left(\xi_{\gamma'\varphi(f(x))=x')}\oplus R_m\right).  
\end{array}
\end{equation*}
Since the decompositions of  both $\mathbb{V}(\Gamma)$ and $\mathbb{V}(\Gamma')$  in $\left(\xi_{\gamma(x)}\oplus R_m\right), x\in V$ and  
$\left(\xi_{\gamma'(x')}\oplus R_m\right),$ $\varphi(x)=x'\in V'$ is a bijection $(\varphi; \Id)$  can be extend to a $\overline{\mathbb{F}}_p$-isomorphism 
$\varphi : \mathbb{V}(\Gamma_m) \longrightarrow \mathbb{V}(\Gamma'_m)$.
  
 \noindent We can also see, considering that the above isomorphism is an isomorphism of abelian groups,   that for all $i\in\mathcal{E}_m$ we have:  ${\displaystyle \varphi(\alpha(i))=\alpha'(i)}$, so that
${\displaystyle  x \in \alpha(i) \Longleftrightarrow \varphi(x)\in \; \alpha'(i)}$ : 
${\displaystyle  \gamma(x)=\gamma'(\varphi(x))}$; so ${\displaystyle  \varphi(i.x)=\varphi(\chi_{\gamma(x)}(i)x)=\chi_{\gamma'(\varphi(x))}(i)\varphi(x)=i.\varphi(x)}$ : $\varphi$ is an isomorphism of $R$-modules;


\item[-] ii)$\Longrightarrow$ i) : suppose now that  $\mathbb{V}(\Gamma_m) \simeq \mathbb{V}(\Gamma'_m)$ as $R_m$-modules, they have the same canonical decomposition, of the form 
$$\mathbb{V}(\Gamma)\simeq \bigoplus_{a \in \mathbb{E}_m}m_a (\xi_a\oplus R_m)=
 m_{is} \left(\xi_{0}\oplus R_m\right)\bigoplus_{e\in I_{\mathcal{E}_m}} 2 \left(\xi_{e}\oplus R_m\right) \bigoplus_{x \in V_{\gamma}} \left(\xi_{\gamma(x)}\oplus R_m\right)$$
and
$$\mathbb{V}(\Gamma')\simeq \bigoplus_{a \in \mathbb{E}_m}m'_a (\xi_a\oplus R_m)=
 m'_{is} \left(\xi_{0}\oplus R_m\right)\bigoplus_{e'\in I'_{\mathcal{E}_m}} 2 \left(\xi_{e}\oplus R_m\right) \bigoplus_{x' \in V'_{\gamma'}} \left(\xi_{\gamma'(x')}\oplus R_m\right).$$
So we have $m_a=m'_a$ for all $a$.
\begin{itemize}
\item[-] Firstly the coefficients of $\xi_0 \oplus R_m$ are the same so $m_{is}=m'_{is}: \Gamma$ and $\Gamma'$ have the same number of isolated vertices;
\item[-] Secondly $\{a: m_a=2\}=\{ a: m'_a=2\}$ so that $ I_{\mathcal{E}_m}= I'_{\mathcal{E}_m}$ and  $\Gamma, \Gamma'$ have the same number of isolated edges;
\item[-] Thirdly $\{a: m_a=1\}=\{ a: m'_a=1\}$ and for all  $x \in V_{\gamma}$ there exists an  unique $x' \in
 V'_{\gamma'}$ such that $\xi_{\gamma(x)}=\xi_{\gamma'(x')}$, that is $\gamma(x)
 =\gamma'(x')$ which   gives a bijection: 
\end{itemize}
 \begin{eqnarray*} 
 \varphi  : V_{\gamma}& \longrightarrow V'_{\gamma'}\\
                               x    & \longmapsto x'
\end{eqnarray*}
The map $\varphi$ preserves the  incidence : if $i$ is  incident to $x$ in $\Gamma_m$, we have $i \in
\gamma(x)$; so, from above  we have: $i \in \gamma'(x')$: $i$ is incident to $x'$ in $\Gamma'_m$ : $\varphi\circ\alpha(i)=\alpha'(i)$.\\

\noindent Finally, we complete $\varphi$ by a bijection between isolated vertices, and a bijection between isolated edges.
\item[-] $ii)\Longleftrightarrow $iii) : from above. \qedhere
\end{itemize}
\end{proof}

\end{document}
